\documentclass[11pt]{article}

\usepackage{geometry}
\geometry{margin=1in}
\usepackage{lmodern}
\usepackage[T1]{fontenc}
\usepackage[section]{placeins}	%Prevents figures from occurring in other sections with the \FloatBarrier command
\usepackage{graphicx}	% lets you put in figures and gives access to an array of figure formatting options
\graphicspath{ {./Figures/} }	%Figures are saved in a folder named "Figures" under the directory of the main document
\usepackage{braket}
\usepackage{amsmath}	% Math libraries
\usepackage{amssymb}	% Math libraries
\usepackage{mathbbol}	% Math libraries
\usepackage{mathtools}	% Math libraries
\usepackage{textcomp}	% Symbols libraries
\usepackage{enumitem}	% Customizes the description environment; lets me align text within a descriptive list (list with no bullet points).
\usepackage{hyperref}	% Reference tool. Lets you dynamically link references to superscripts in the document text
\usepackage{bm}	% this one is important somehow....
\usepackage{color}	 % allows you to use text color commands - good for edits
\usepackage[capposition=bottom]{floatrow}	% Figure captions positioned below the figure
\usepackage{setspace}	% Sets vertical spacing in text
\usepackage{achemso}	% ACS citation format package
\setkeys{acs}{articletitle=true}	% Boolean statement that controls the addition of article titles to citations in the bibliography
\usepackage{tocloft}	% Adds dot leaders
\usepackage{blindtext}

\title{Molecular Mechanisms Driving Polysaccharide Enrichment in Sea Spray Aerosol Proxy Systems}
\date{\vspace{-8ex}}

\begin{document}

\maketitle

Sea spray aerosol (SSA) is one of the most abundant sources of natural aerosol with significant but poorly understood impacts on Earth's radiative energy budget.\cite{deleeuwProductionFluxSea2011,ipcc2014ClimateChange20132014} Organic compounds produced by marine biota partition to the air-sea interface, known as the sea surface microlayer (SSML), where the organic matter is transferred to SSA upon wave breaking.\cite{quinnChemistryRelatedProperties2015} These organic molecules form a film around the aqueous salty core of the particle.\cite{tervahattuNewEvidenceOrganic2002} Although sugars are generally regarded as water soluble and surface inactive, sugars have been measured with high enrichment factors in SSA as compared to seawater.\cite{russellCarbohydratelikeCompositionSubmicron2010,jayarathneEnrichmentSaccharidesDivalent2016} To determine the molecular mechanisms behind sugar enrichment in SSA, a proxy SSML system was developed. SSA is coated in highly surface-active aliphatic compounds, with palmitic acid being the most abundant surfactant in SSA and the SSML.\cite{cochranAnalysisOrganicAnionic2016} Na$^+$, Ca$^{2+}$, and Cl$^-$ comprise three of the most prevalent ions in seawater.\cite{milleroCompositionStandardSeawater2008a,milleroChemicalOceanography2013} Alginate serves as a marine-relevant model polysaccharide that readily forms gels in the presence of Ca$^{2+}$.\cite{fangMultipleStepsCritical2007,liCrossLinkedPolysaccharideAssemblies2013,sunMolecularDynamicsSimulations2014,allerSizeresolvedCharacterizationPolysaccharidic2017} The polysaccharide was enriched in laboratory-generated SSA, and its enrichment factor was higher in the presence of CaCl$_2$.\cite{haseneczSeaSprayAerosol2018} Consequently, the molecular organization of palmitic acid monolayers on the surface of NaCl aqueous solutions was monitored for changes upon introducing alginate and alginate-Ca$^{2+}$ gels into the system. Changes in film molecular organization yield information regarding the potential mechanisms of polysaccharide enrichment in SSA. \par

Multiple surface-sensitive techniques were used to study changes in palmitic acid film organization upon the addition of alginate and CaCl$_2$. Pendant drop surface tensiometry measured equilibrium surface tension of palmitic-acid-coated salty water droplets with alginate and alginate-Ca$^{2+}$ gels within the droplet. Infrared reflection-absorption spectroscopy (IRRAS) and vibrational sum frequency generation spectroscopy (VSFG) probed interfacial molecular organization via functional group vibrational signatures. More specifically, IRRAS was used to determine palmitic acid lattice packing at the air-water interface as a function of Ca$^{2+}$ and alginate presence.\cite{adamsSodiumCarboxylateContact2017} VSFG was used to determine palmitic acid alkyl chain orientation.\cite{adamsSodiumCarboxylateContact2017} The surface propensity of alginate was also confirmed by IRRAS and VSFG spectra. Lastly, Brewster angle microscopy (BAM) was used to visualize changes in film morphology after introducing alginate and Ca$^{2+}$ to the system.\cite{ruddThermodynamicNonequilibriumStability2018} Alginate was surface active, as demonstrated by the small surface tension depression of the fatty-acid-coated salty water droplet. Additionally, alginate disordered the film and increased the lipid tilt angle, indicating that the polysaccharide interacts with the palmitic acid monolayer. BAM images revealed that a small amount of alginate intercalated between surfactants in the monolayer; however, the majority of the alginate likely adsorbed to the palmitic acid head groups immersed in the aqueous phase of the system. Adding CaCl$_2$ produced greater surface tension depression and more gel particles visible at the surface. Larger palmitic acid tilt angles and significant monolayer disorder were also observed. Alginate-Ca$^{2+}$ gels are likely highly surface active and intercalate into the monolayer due to their more compact structures. Alginate in the absence of calcium does not form such tight gels, so it is generally too large to intercalate between the monolayer lipids and simply adsorbs to the head groups of the film. Therefore, it is expected that intermolecular interactions between the polysaccharide gels and fatty acid films and gel intercalation into the SSML likely account for the high sugar enrichment in SSA. \par

%Multiple surface-sensitive techniques were used to study changes in palmitic acid film organization upon the addition of alginate and CaCl$_2$. Pendant drop surface tensiometry measured equilibrium surface tension of palmitic-acid-coated salty water droplets with alginate and alginate-Ca$^{2+}$ gels within the droplet. Infrared reflection-absorption spectroscopy (IRRAS) and vibrational sum frequency generation spectroscopy (VSFG) probed interfacial molecular organization via functional group vibrational signatures. More specifically, IRRAS was used to measure the C-D scissoring modes of deuterated palmitic acid to determine its lattice packing at the air-water interface as a function of Ca$^{2+}$ and alginate presence.\cite{adamsSodiumCarboxylateContact2017} VSFG was used to determine palmitic acid alkyl chain orientation via the C-D stretching modes.\cite{adamsSodiumCarboxylateContact2017} The surface propensity of alginate was also confirmed by observing its C-H stretching modes in both IRRAS and VSFG spectra. Lastly, Brewster angle microscopy (BAM) was used to visualize changes in film morphology after introducing alginate and Ca$^{2+}$ to the system.\cite{ruddThermodynamicNonequilibriumStability2018} Increasing concentrations of alginate in NaCl solutions yielded $\sim$5 mN/m surface tension depression, some monolayer disorder, and greater palmitic acid alkyl chain tilt angles as compared to the fatty acid monolayer alone. BAM images revealed that some alginate intercalated between surfactants in the monolayer; however, the majority of the alginate likely adsorbed to the palmitic acid head groups. Adding CaCl$_2$ produced greater surface tension depression and more gel particles visible at the surface. Larger palmitic acid tilt angles and significant monolayer disorder were also observed. Alginate-Ca$^{2+}$ gels are likely highly surface active and intercalate into the monolayer due to their more compact structures. Alginate in the absence of calcium does not form such tight gels, so it is generally too large to intercalate between the monolayer lipids and simply adsorbs to the head groups of the film. Therefore, it is expected that intermolecular interactions between the polysaccharide gels and fatty acid films and gel intercalation into the SSML likely account for the high sugar enrichment in SSA. \par

%Multiple surface-sensitive techniques will be used to study the changes in palmitic acid film organization upon the addition of alginate and CaCl$_2$. Surface tensiometry methods, such as a pendant drop tensiometer, measure surface tension and other parameters necessary for evaluating interfacial thermodynamics. The equilibrium surface tension of palmitic acid-coated salty water droplets containing variable amounts of alginate will be measured to understand how alginate impacts droplet surface tension. Surface-sensitive spectroscopy, such as infrared reflection-absorption spectroscopy (IRRAS) and sum frequency generation spectroscopy (SFG), probe interfacial molecular organization via functional group vibrational signatures. Deuterated palmitic acid will be purchased to enable lattice packing analysis of the C-D scissoring mode via IRRAS.\cite{adamsSodiumCarboxylateContact2017} Secondly, palmitic acid alkyl chain orientation will be measured via the C-D stretching modes with SFG to determine how alginate and Ca$_2$ affect the monolayer organization. IRRAS and SFG will also be used to measure the surface propensity of alginate via its C-H stretching modes. Because IRRAS probes a greater interfacial depth than SFG, more information about alginate interfacial presence can be gleaned from both methods.\\

%Motivation: Saccharides are enriched in SSA. Alginate, a polysaccharide, is enriched in SSA generated in a mini-MART mesocosm experiment, and the enrichment factor is higher upon adding calcium chloride and reef salt.

%Question: How does alginate interact with palmitic acid monolayers in the presence and absence of CaCl$_2$?

%System: Alginate (1 ppm, 5 ppm, 10 ppm, 25 ppm, 50 ppm), palmitic acid, 470 mM NaCl, 10 mM CaCl$_2$, pH 8.2

%Measurements: Equilibrium surface tension (pendant drop), images of monolayer reorganization (BAM), C-D scissoring mode for lattice packing (IRRAS), C-D stretching modes for alkyl chain orientation (SFG, ppp \& ssp), C-H stretching modes for alginate (IRRAS \& SFG) \\

\bibliography{HayesForum2019}

\end{document}